\documentclass[11pt]{article}

\usepackage[margin=1in]{geometry}
\usepackage{amsmath,amssymb}
\usepackage{booktabs}
\usepackage{graphicx}
\usepackage{xcolor}
\usepackage{listings}
\usepackage{lmodern}
\usepackage[T1]{fontenc}
\usepackage{microtype}
\usepackage[hidelinks]{hyperref}
\usepackage[capitalise]{cleveref}
\usepackage{natbib}

\lstset{
  basicstyle=\ttfamily\small,
  breaklines=true,
  columns=fullflexible,
  frame=single,
  framesep=4pt,
  xleftmargin=6pt,
  literate={~}{{\textasciitilde}}1
}

\newcommand{\pass}{\textsc{pass}}
\newcommand{\usable}{\textsc{usable}}
\newcommand{\vuln}{\textsc{vulnerable}}
\newcommand{\dfaeq}{\textsc{dfa-eq}}
\newcommand{\exact}{\textsc{exact}}

\title{\textbf{Passing Is Not Shipping}\\[0.4em]
\large Correctness, Semantic Equivalence and ReDoS Safety in\\
LLM-Generated Regular Expressions, and the Limits of Measuring Them}

\author{Plicara Labs\\\texttt{info@plicara.ai}}

\date{August 2026}

% Numbers that would otherwise be transcribed by hand. Generated by
% paper/make_tables.py from committed results; CI fails if they drift.
% generated by make_tables.py -- do not edit
\newcommand{\anchoredgold}{13.4}
\newcommand{\anchoredmodels}{9.8}
\newcommand{\anchoredmodelsci}{1.0}
\newcommand{\anchoredmodelsn}{3{,}613}
\newcommand{\anchoredp}{0.20}
\newcommand{\anchoredprod}{8.9}
\newcommand{\auditcomplement}{85}
\newcommand{\bandwidth}{6.7}
\newcommand{\bhcount}{8}
\newcommand{\cheapusable}{19.8}
\newcommand{\commontasks}{421}
\newcommand{\correctsamples}{5{,}269}
\newcommand{\correctvulnsamples}{390}
\newcommand{\corrgold}{15.5}
\newcommand{\corrlib}{21.9}
\newcommand{\corrmodels}{13.4}
\newcommand{\corrprod}{11.1}
\newcommand{\corrso}{18.7}
\newcommand{\costgap}{1.1}
\newcommand{\dearusable}{20.8}
\newcommand{\discriminative}{162}
\newcommand{\dropanchoredpct}{23.7}
\newcommand{\droprange}{\texttt{cpan} at 7.7\% down to \texttt{pypi} at 0.3\%}
\newcommand{\easyhard}{-3.6}
\newcommand{\easyhardhi}{+3.8}
\newcommand{\easyhardlo}{-12.2}
\newcommand{\easyvuln}{7.1}
\newcommand{\equivloss}{18.9}
\newcommand{\equivshare}{87}
\newcommand{\exacthi}{5.0}
\newcommand{\exactlo}{1.8}
\newcommand{\examplediff}{+1.2}
\newcommand{\examplehi}{+4.8}
\newcommand{\examplelo}{-2.6}
\newcommand{\gapclosed}{0.4}
\newcommand{\gapcorr}{10.8}
\newcommand{\gapraw}{11.2}
\newcommand{\hardvuln}{10.7}
\newcommand{\holmcount}{6}
\newcommand{\identicalpass}{56}
\newcommand{\identicalusable}{62}
\newcommand{\identicalvuln}{78}
\newcommand{\keptanchoredpct}{17.2}
\newcommand{\libdropped}{392}
\newcommand{\libdroppedpct}{10.2}
\newcommand{\libpool}{3{,}838}
\newcommand{\linearonly}{20{,}658}
\newcommand{\linearpatterns}{23{,}941}
\newcommand{\linearregistries}{\texttt{crates.io} and \texttt{godoc}}
\newcommand{\modelclusterhi}{12.4}
\newcommand{\modelclusterlo}{7.5}
\newcommand{\modelprodcorr}{2.2}
\newcommand{\modelprodraw}{0.9}
\newcommand{\passhi}{46.5}
\newcommand{\passlo}{38.0}
\newcommand{\proddropped}{26{,}900}
\newcommand{\proddroppedpct}{5.0}
\newcommand{\prodpool}{537{,}804}
\newcommand{\readjacent}{154}
\newcommand{\readjacentpct}{40}
\newcommand{\recallrange}{67--93}
\newcommand{\recovered}{17{,}151}
\newcommand{\recoveredpct}{64}
\newcommand{\relen}{28}
\newcommand{\renested}{236}
\newcommand{\reqgrouppct}{34.7}
\newcommand{\requant}{3}
\newcommand{\resolvedpass}{16}
\newcommand{\resolvedusable}{9}
\newcommand{\resolvedvuln}{9}
\newcommand{\robboth}{8.2}
\newcommand{\robbt}{8.3}
\newcommand{\robnorm}{9.1}
\newcommand{\robrange}{8.2--9.1}
\newcommand{\safetyloss}{2.9}
\newcommand{\sodropped}{56{,}572}
\newcommand{\sodroppedpct}{11.4}
\newcommand{\sopool}{495{,}135}
\newcommand{\sradjacent}{355}
\newcommand{\sradjacentpct}{64}
\newcommand{\srlen}{33}
\newcommand{\srnested}{202}
\newcommand{\srqgrouppct}{18.6}
\newcommand{\srquant}{4}
\newcommand{\srvulngivencorrect}{16.5}
\newcommand{\topunrecovered}{malformed or truncated by extraction}
\newcommand{\topunrecoveredn}{4{,}588}
\newcommand{\unadjustedcount}{8}
\newcommand{\undecshare}{48}
\newcommand{\undecsplit}{451 \textsc{unsup}, 20 \textsc{undec}}
\newcommand{\undecsupported}{471}
\newcommand{\undecusable}{974}
\newcommand{\vgchi}{10.0}
\newcommand{\vgclo}{5.6}
\newcommand{\vulngivencorrect}{7.4}


\begin{document}
\maketitle

\begin{abstract}
\noindent
We evaluate eleven contemporary language models on natural-language-to-regex
generation over 450 tasks with three samples each ($14{,}850$ requests,
\$10.95), scoring every answer on three axes: whether it satisfies the task's
tests, whether it denotes the same language as a human reference, and whether
it is vulnerable to regular-expression denial of service (ReDoS).

Our headline result is a negative one about a composite metric we built from
the corpus's own metric suite and then had to abandon. A composite requiring correctness,
safety and semantic equivalence appears to show that ``passing is twice as
easy as shipping'': models pass 38.0--46.5\% of tasks but satisfy all three
criteria on only 17.1--23.8\%. Decomposing that gap shows that
\textbf{\equivshare{}\% of it is contributed by the semantic-equivalence term}, and a
manual audit shows that term is itself dominated by reference-set error and
task underspecification rather than model behaviour. Using only the two
criteria that never consult a human answer key, the gap is
\textbf{\safetyloss{} percentage points}: just \vulngivencorrect{}\% of functionally correct patterns
are ReDoS-vulnerable.

That figure is itself informative by contrast. Joint correctness-and-security
benchmarks in general code generation report far larger penalties. BaxBench
finds roughly half of correct backends exploitable. SecureAgentBench reports
15.2\% correct-and-secure for its best agent. The regex domain does \emph{not}
reproduce that pattern, and we report the difference rather than the
similarity. Separately, models are safer than the corpus's own human-written
reference patterns --- 9.0\% pooled against 13.6\%, and eight of the eleven
separate under an exact McNemar test on the paired tasks. We initially read
that as evidence that ReDoS is endemic to human practice. Extending the
identical screen to five further populations, including half a million
regular expressions extracted from shipped packages, refutes it and replaces
it with a sharper ordering. Controlling for task mix by restricting every
population, models included, to anchored patterns: those \emph{written to be
read} are ReDoS-prone (reusable library entries 20.1\%, Stack Overflow
answers 17.3\%, this corpus's reference answers \anchoredgold{}\%), while those
\emph{executed} in shipped packages sit at \anchoredprod{}\%, statistically
indistinguishable from the models at \anchoredmodels{}\%. \emph{Vulnerability tracks
whether a pattern was ever run, not whether a human or a model wrote it.} We
calibrate the screen against an independent detector and a dynamic oracle to
show the ordering is not an artifact of differential miss rates.

The audit behind this is the paper's methodological contribution. Metrics
that compare generated output against a human answer key inherit that key's
error rate, a dependency documented for image and text benchmarks
\citep{northcutt2021labelerrors} but not, to our knowledge, quantified for
this one. In a random sample of disagreements where a model passed every
test yet was scored as semantically different, the model was clearly at fault
in roughly 15\% of cases. The rest split between demonstrably incorrect
reference patterns (including a reference character class containing
a literal \texttt{|}) and prompts that never specified the property in
dispute. Because that unreliable term supplies \equivshare{}\% of our composite's
headline gap, the composite was measuring its own weakest component. We also show that independent binomial intervals, the default in this
literature, are the wrong inferential tool when all systems answer identical
items, and that a paired bootstrap resolves nine times as
many pairwise comparisons on the same data.

Replicating the two reference-independent criteria on a second corpus
\citep{ye2020structuredregex} shows the qualitative finding holds and the
number does not: the same eleven models leave \textbf{\srvulngivencorrect{}\%} of their
correct patterns ReDoS-vulnerable there against \vulngivencorrect{}\% here, on tasks twenty
points \emph{easier}. The correctness-to-security penalty is a property of
the corpus, not of the models.

We therefore decline to publish a ranking. We release all raw generations,
scoring code and adjudications, and recommend that evaluations separate
reference-dependent from reference-independent metrics, assign them different
evidentiary weight, and treat any effect size as corpus-local until
replicated.
\end{abstract}

\vspace{0.5em}
\noindent\textbf{Artifacts.} All model outputs, scoring code, per-case
adjudications and analysis scripts are available at
\url{https://github.com/plicara/regexeval-2026}~\citep{regexeval-2026}. Every number in
this paper recomputes offline from committed data with no API access.

\section{Introduction}

Regular expressions make an unusually good probe for language-model
evaluation. The task is compact. The output runs. Correctness is partly
decidable, which is rare. And a syntactically valid, test-passing answer can
still be a security defect, which is what drew us to the task in the first
place.

The standard protocol for this task scores a candidate against a reference
by DFA equivalence~\citep{locascio2016neural,kushman2013semantic}, sometimes
alongside test-based correctness. Neither criterion notices catastrophic
backtracking. A pattern can satisfy every provided example, denote exactly the
reference language, and still admit inputs on which matching takes time
exponential in the input. That is the ReDoS vulnerability
class~\citep{davis2018redos}, a denial-of-service vector in deployed systems
and common enough in human-written code to have prompted ecosystem-scale
study.

\citet{siddiq2024regexeval} closed that gap. Their Re(gEx$|$DoS)Eval
introduces the corpus we use, defines the four metrics we use
($\pass{}@k$, $\vuln{}@k$, $\dfaeq{}@k$ and exact match), scores correctness
and security jointly on them, and evaluates T5, Phi-1.5 and GPT-3.5-Turbo.
The instrument in this paper is theirs. We say so plainly here because an
earlier draft of this work did not, and because it determines what is left
for us to contribute: not the measuring apparatus, but what happens when it
is turned on the current model population, decomposed, and then turned back
on itself.

\subsection{Contributions}

\begin{enumerate}
  \item \textbf{A decomposition of the joint metric}
  (\cref{sec:decomp}). The composite is the natural way to report
  correctness and security together and it is what the joint-benchmark
  literature reports (\cref{sec:jointprior}). We show that on this corpus the
  composite is dominated by its equivalence term, which supplies \equivshare{}\% of the
  headline gap, and that the term is the least trustworthy of the three. The
  contribution is the finding that the aggregate is uninformative here, not
  the aggregate.

  \item \textbf{A quantified audit of reference-standard error}
  (\cref{sec:goldstandard}). We sample disagreements and adjudicate them,
  finding that the majority of apparent semantic errors are not model errors.
  \citet{northcutt2021labelerrors} establish that test-set label error is
  pervasive and can reverse benchmark orderings; we measure it for a corpus
  whose labels are regular expressions, where a label error is a formal
  object that can be exhibited and argued with rather than a disputed
  annotation. \citet{siddiq2024regexeval} report their metrics without
  auditing the reference set they compare against, which is the standard
  practice we are questioning rather than a defect specific to them.

  \item \textbf{A human baseline for ReDoS across six populations}
  (\cref{sec:humanbaseline}). We apply the identical safety screen to the
  corpus's own reference patterns, a second NL-to-regex gold set, a
  grammar-generated control, a reusable-pattern library, Stack Overflow, and
  half a million regular expressions from shipped packages. The resulting
  ordering is monotone in whether a pattern is executed rather than in who
  wrote it, which identifies the mechanism behind the reference set's
  unreliability instead of only observing it. Screening costs no inference,
  so this generalises where the model-side results cannot.

  \item \textbf{An updated model population.}
  \citet{siddiq2024regexeval} evaluate T5, Phi-1.5 and GPT-3.5-Turbo. We
  evaluate eleven models released in the intervening period under pinned
  serving providers (\cref{sec:serving}), and find the population compressed
  into an eight-point band that a hundredfold price range does not order.

  \item \textbf{A second-corpus replication} (\cref{sec:replication}). We
  re-run the two reference-independent criteria on
  StructuredRegex~\citep{ye2020structuredregex} with the same models and
  configuration. The finding replicates in kind and not in size: \srvulngivencorrect{}\%
  against \vulngivencorrect{}\%, on easier tasks, with the model ordering not preserved. This
  is the direct test of \cref{sec:crossbench}'s domain-dependence claim,
  which otherwise rests on comparison across languages and vulnerability
  classes.

  \item \textbf{Full artifact release} including raw generations, which
  permits any of our scoring decisions to be overturned without re-running
  inference.
\end{enumerate}

\subsection{What this paper does not claim}

We do not claim the metrics. $\pass{}@k$ is \citet{chen2021codex};
$\vuln{}@k$, $\dfaeq{}@k$ and exact match on this corpus are
\citet{siddiq2024regexeval}. Joint correctness-and-security scoring is
established across several domains (\cref{sec:jointprior}) and on this one.
The \usable{} composite in \cref{sec:metrics} is a conjunction of their three
metrics, and the argument of this paper is that it should not be used.

We do not claim the paired bootstrap of \cref{sec:paired}. It is
\citet{bergkirkpatrick2012empirical}, following \citet{koehn2004bootstrap},
and its use is recommended practice~\citep{dror2018hitchhiker}. What we
report is how much it changes on this data: nine resolvable pairwise
comparisons against one under independent intervals. That is an observation
about the corpus, not a method.

We do not publish a ranking of the eleven models. \cref{sec:power} shows
that 62\% of tasks yield identical outcomes across all eleven systems, and
that only nine of fifty-five pairwise comparisons resolve at 95\% confidence
even under the correct paired test. Reporting an ordinal leaderboard from
this data would be unsupported by it.

We do not claim novelty for DFA-equivalence checking, which is standard
in this literature~\citep{locascio2016neural}, nor for ReDoS detection,
which has a long static-analysis
lineage~\citep{kirrage2013static,wustholz2017static,sok2025redos}, nor for
the scoring engine, which is prior work by the present authors released
separately~\citep{regexbench} and used here as a pinned dependency.

\section{Background and Related Work}

\paragraph{NL-to-regex generation.}
The task was established by \citet{kushman2013semantic} (KB13) and extended
by \citet{locascio2016neural} (NL-RX), both of which adopt DFA equivalence
against a reference as the correctness criterion. That is the right choice in principle. $\texttt{[0-9]+}$ and
$\texttt{[0-9][0-9]*}$ denote the same language and differ as strings, so
text comparison would mark one of them wrong. But the criterion inherits any
defect in the reference set, which is the dependency this paper quantifies.

\paragraph{Re(gEx$|$DoS)Eval.}
\citet{siddiq2024regexeval} supply both the corpus and the metric suite used
here. They collect 762 tasks from real user posts, each with a description,
positive and negative test strings and a human reference pattern, and score
generations on $\pass{}@k$, $\vuln{}@k$, $\dfaeq{}@k$ and exact match,
reporting which model produces regexes that are correct \emph{and} secure.
Their evaluation covers T5, Phi-1.5 and GPT-3.5-Turbo. Everything this paper
measures is measured with their instrument. Our departures are three: we run
it on eleven models released since, we decompose the joint metric instead of
reporting it whole (\cref{sec:decomp}), and we turn the safety screen on
their reference patterns (\cref{sec:humanbaseline}), which they do not do.
The companion study \citep{siddiq2024redos} examines ReDoS in LLM-generated
patterns alongside developer-forum discussion and reports a skew toward
polynomial over exponential vulnerability; we compare against that finding
in \cref{sec:humanbaseline}.

\paragraph{ReDoS.}
Catastrophic backtracking arises when a backtracking matcher can decompose a
prefix in exponentially many ways, canonically via a quantifier applied to a
quantified group, as in \texttt{(a+)+}. Detecting it statically is a
developed area: \citet{kirrage2013static} give an analysis for exponential
blow-up by search over the pattern's derivation tree,
\citet{wustholz2017static} extend detection to programs that construct
patterns dynamically, and \citet{sok2025redos} survey the detector and
mitigation literature. The scoring engine we depend on
\citep{regexbench} sits in that lineage and combines a structural screen with
witness-string execution. On prevalence, \citet{davis2018redos} document
ReDoS at ecosystem scale and \citet{davis2019linguafranca} study regex
portability and re-use across languages.

\paragraph{Security of generated code.}
The observation that models emit working but unsafe code predates the joint
benchmarks below. \citet{pearce2022asleep} generate 1{,}689 programs across
89 security-relevant scenarios and find roughly 40\% vulnerable, and
\citet{perry2023users} run a controlled user study showing that participants
with an AI assistant wrote less secure code while believing the opposite.
Both motivate scoring security separately from correctness rather than
assuming the first implies the second.

\subsection{Joint correctness-and-security benchmarks}
\label{sec:jointprior}

Evaluating functional correctness and security together is an active area, and
this paper is not its first instance. \citet{peng2025cweval} introduce CWEval,
scoring 119 security-critical tasks across 31 CWEs and five languages with
outcome-driven oracles for both properties, and report a substantial
population of functionally correct but insecure generations.
\citet{vero2025baxbench} present BaxBench, 392 backend tasks with expert
exploits, and find that roughly half of functionally correct solutions are
exploitable. \citet{chen2025secureagentbench} evaluate coding \emph{agents} on
105 repository-level tasks, reporting 15.2\% correct-and-secure for the best
agent and 9.2\% on average. \citet{patir2025dualgauge} automate the joint
pipeline in DualGauge over 154 tasks and 10 models, observing secure-pass@1
below 12\% while functional pass@1 exceeds 50\%.

The joint framing is therefore established both across domains and, by
\citet{siddiq2024regexeval}, on this one. What differs here is the shape of
the instrument. The artifact is a single expression instead of a program. The
vulnerability class is ReDoS instead of a CWE taxonomy. Correctness is partly
decidable, not established by tests alone. And a human reference exists for
every task, which the general-code benchmarks do not have: they execute tests
and exploits rather than compare against a gold artifact. That difference
cuts both ways. It makes \cref{sec:humanbaseline} possible, and it is the
sole reason \cref{sec:goldstandard} is necessary. The reference-free
benchmarks are immune to the failure mode we spend half this paper
characterising.

Our numbers turned out not to match theirs. \cref{sec:crossbench} works
through the difference.

\paragraph{Sampling-based metrics.}
We use the unbiased $\pass{}@k$ estimator of \citet{chen2021codex}, with the
degeneracy noted in \cref{sec:metrics}.

\paragraph{Significance testing.}
Comparing systems that answer identical items calls for a paired procedure.
The paired bootstrap we use is \citet{bergkirkpatrick2012empirical},
following the bootstrap-resampling protocol of \citet{koehn2004bootstrap} for
machine translation, and its use over unpaired intervals is standard
advice~\citep{dror2018hitchhiker}; the discrete analogue for paired binary
outcomes is \citet{mcnemar1947note}. \cref{sec:paired} reports what the
choice costs on this corpus rather than proposing anything new.

\paragraph{Contamination.}
This corpus was published in 2024 and its tasks come from public forum posts,
so it is plausibly inside the training data of every model we evaluate.
\citet{sainz2023contamination} argue that contamination must be measured per
benchmark rather than assumed away. We cannot measure it here and say so in
\cref{sec:threats}.

\paragraph{Serving-layer variance.}
Hosted inference routers may serve the same model slug from multiple
providers at differing numerical precision. \citet{openrouterpinning} report
that 31 of 32 surveyed AI-safety codebases using such a router did not pin
the serving provider, making their results a measurement of the router's
routing policy in addition to the model. \cref{sec:serving} describes the
discipline we adopt in response.

\section{Method}

\subsection{Task and prompt}

Each task supplies a natural-language description. The model returns one
regular expression. The prompt is fixed across all models and all tasks:

\begin{quote}\small
\textbf{system:} You translate a natural-language description into a single
Python \texttt{re}-compatible regular expression. Reply with ONLY the pattern
in one fenced code block, no explanation.\\[0.3em]
\textbf{user:} \textit{\{task description, verbatim\}}\\[0.3em]
Reply with only the regular expression pattern in a single fenced code block.
\end{quote}

We do not tune this prompt. Prompt optimisation would raise all scores and
destroy comparability with prior work on the same corpus. \cref{sec:threats}
notes the corresponding threat.

\paragraph{Extraction.} We take the final fenced code block, else the whole
trimmed reply, then strip at most one layer of host-language quoting
(\texttt{r'\ldots'}, \texttt{'\ldots'}, \texttt{"\ldots"}, backticks,
\texttt{/\ldots/flags}). Scored literally, \texttt{r'\textbackslash d+\$'} is
a pattern matching the letter \texttt{r}, and would fail for reasons
unrelated to regex competence. This affected 7--18 of 1{,}350 responses per
model, under 1.4\% throughout, and changes no conclusion. We release both the
normalised and the unnormalised scores.

\subsection{Corpus}

We use Re(gEx$|$DoS)Eval~\citep{siddiq2024regexeval}, 762 tasks drawn from real user
requests, each with a description, positive and negative test strings, and a
human-written reference pattern. We evaluate 450 tasks, selected at uniform
stride across the corpus rather than as a prefix (the corpus is
difficulty-ordered at its head). The subset size was budget-determined.

Match semantics are set per corpus. Under search semantics 761 of the 762
Re(gEx$|$DoS)Eval references satisfy their own tests, against 94\% under
full-match. Choose wrong and 46 gold patterns are scored as failures. We
verified the loaded configuration by scoring every reference against its own
tests before the run (\cref{sec:goldstandard}).

The single exception is worth its own sentence, because it is not a mismatch.
\texttt{regexeval/1660}'s reference,
\lstinline!^((\.)?([a-zA-Z0-9_-]?)(\.)?([a-zA-Z0-9_-]?)(\.)?)+$!, does not
accept a string it should reject: it fails to \emph{finish} on one, exceeding
the scorer's one-second match timeout, which is counted as a false positive.
The safety screen independently flags the same pattern as exponential. The one
reference in the corpus that fails its own tests fails because it is
ReDoS-vulnerable.

\subsection{Models and serving discipline}
\label{sec:serving}

Eleven models spanning frontier and open-weights systems across a $98\times$
price range. Every request pins a single named provider with fallback
disabled, so a request is either served by the named endpoint or fails
visibly. Substitution failures did occur and are reported as failures rather
than silently re-routed (\cref{app:failures}).

For open-weights models the pin includes quantisation where the router
discloses it: GLM-5.2 and DeepSeek-V4-Flash are served at 4-bit by some
providers and 8-bit by others, and we pin 8-bit. One pin went stale within an
hour of being set (the provider ceased serving that model), so pins are
re-verified immediately before each run.

The pin is a request, not a guarantee. So we record which provider actually
served each response and check afterwards that it matches. All eleven models
came back served entirely by the endpoint they were pinned to.

\subsection{Metrics}
\label{sec:metrics}

Let $T=(d, P, N, r)$ be a task with description $d$, positive examples $P$,
negative examples $N$ and reference pattern $r$. Let $p$ be a candidate.

\paragraph{Correctness.}
\[
\mathrm{correct}(p,T) \;=\; \big[\forall x \in P:\ \mathrm{match}(p,x)\big]
\;\wedge\;
\big[\forall y \in N:\ \neg\,\mathrm{match}(p,y)\big]
\]
evaluated with CPython's \texttt{re} under the corpus's declared semantics.
This is \emph{reference-independent}.

\paragraph{Semantic equivalence.}
Both $p$ and $r$ are compiled to finite automata and compared via
\texttt{dk.brics.automaton}~\citep{brics}, yielding a
verdict in $\{\textsc{eq}, \textsc{diff}, \textsc{undec}, \textsc{unsup}\}$.
When the verdict is \textsc{diff} the engine returns a \emph{witness}: a
shortest string in the symmetric difference, which renders the judgment
falsifiable by inspection.

Equivalence is decidable only on the regular subset. Backreferences make the
language non-regular, and equivalence then has no algorithmic answer at all.
Those cases return \textsc{undec}. Between 78 and 133 of each model's
scored tasks land there.
We therefore report two figures:
\[
\dfaeq{} = \Pr[\textsc{eq}], \qquad
\dfaeq{}_{\text{dec}} = \Pr[\textsc{eq} \mid \text{verdict} \neq \textsc{undec}].
\]
The first counts undecidable comparisons as failures. It is a lower bound
and cannot flatter. The second isolates the model from the engine's reach.
Report only the second and you inflate. Report only the first and you charge
the model for a theorem. This metric is \emph{reference-dependent}, which
\cref{sec:goldstandard} shows to be decisive.

\paragraph{ReDoS safety.}
$\mathrm{vuln}(p) = [\mathrm{screen}(p) \neq \textsc{safe}]$, where
$\mathrm{screen}$ performs structural analysis for known-vulnerable shapes
followed by an empirical pass against attack strings under timeout. The
structural pass covers three of the five families catalogued
by~\citet{siddiq2024redos}. So \vuln{} is a \emph{lower bound}, and
\textsc{safe} denotes a screening outcome, not a proof. This metric is
\emph{reference-independent}. \cref{sec:calibration} measures how loose the
bound is, separately for each population we compare, because a lower bound
that is looser in one population than another would produce an ordering by
itself.

Where the screen reports a degree, exponential against polynomial, the
structural pass infers it from the matched shape and the empirical pass infers
it from which probe length first exceeded the timeout. The second is a
threshold rather than a measurement of growth order, and we discount the
degree split accordingly (\cref{sec:calibration}).

\paragraph{Composite.}
\[
\mathrm{usable}(p,T) = \mathrm{correct}(p,T)
  \;\wedge\; \neg\,\mathrm{vuln}(p)
  \;\wedge\; \big[\text{verdict}(p,r) \neq \textsc{diff}\big]
\]
The third conjunct excludes only \emph{proven} difference. Undecidable and
unsupported verdicts do not disqualify a pattern. \usable{} is
reference-dependent through that conjunct, and inherits the error
characterised in \cref{sec:goldstandard}.

\paragraph{Sampling.}
With $n$ samples per task of which $c$ satisfy a predicate, we use the
unbiased estimator of~\citet{chen2021codex}:
\[
\widehat{\mathrm{metric}}@k \;=\;
\mathbb{E}_{T}\left[\,1 - \frac{\binom{n-c}{k}}{\binom{n}{k}}\right].
\]
We draw $n=3$ and report $k=3$. At $n=k$ this estimator degenerates to $\mathbf{1}[c \geq 1]$ and
buys no variance reduction over the naive any-of-$n$ statistic. We report it as
$@3$ for comparability, not because the estimator is doing any work at this
setting.
$n>k$ would be required for the estimator's variance properties to apply,
and is left to future work with a larger budget.

The estimator is defined for $n \geq k$, and a handful of tasks returned fewer
samples than that after a refusal or a spending limit. We exclude those tasks
from the affected model's @$k$ estimate rather than score them on what
arrived, and report the count in \cref{tab:at1}. \cref{app:hazards} records
why this is not a cosmetic choice.

\subsection{Sampling configuration}
\label{sec:sampling}

\paragraph{Reasoning disabled.} Four of the eleven models emit no hidden
reasoning tokens. The other seven do. Evaluating them together with reasoning enabled
would confound model identity with inference-time compute budget. All
requests therefore disable reasoning, and reasoning-token counts are recorded
per response to verify compliance (observed: zero throughout).

\paragraph{Temperature unset.} We set
\texttt{require\_parameters=true}, which restricts routing to providers that
honour submitted sampling parameters rather than silently discarding them.
Six of the eleven models reject \texttt{temperature} outright. Together, the
two settings make the request unroutable. We therefore omit \texttt{temperature} entirely and
each model samples at its provider default. That keeps the non-silent-discard guarantee and costs us controlled
sampling. We think it is the right trade and note it as a limitation.

Since $k>1$ is only informative under output diversity, we verified diversity
directly: repeated queries on a fixed task yielded 3, 3, 3 and 2 distinct
patterns for \texttt{claude-opus-5}, \texttt{gpt-5.6-sol}, \texttt{kimi-k3}
and \texttt{glm-5.2} respectively. The analysis pipeline also warns
if $>90\%$ of a model's tasks return identical samples.

\subsection{Controls}

Three synthetic responses go through the identical scoring path on every run.
The task's own reference must come back correct and usable. The pattern
\texttt{z\{5\}} must fail every correctness check. \texttt{(a+)+b} must be
flagged vulnerable. If any control misbehaves we throw the run away instead of
reporting it.

The failure this catches is a scorer that quietly returns nulls, which looks
exactly like every model performing badly. All controls behaved on all eleven
models.

\section{Results}
\label{sec:results}

\subsection{Main results}

\begin{table}[htbp]
\centering
\footnotesize
\setlength{\tabcolsep}{3pt}
\input{tab_main}
\caption{Primary results, $n=k=3$. All figures percentages except $n$ (tasks
entering the @3 estimate) and \emph{undec.} (tasks whose equivalence is
undecidable). Tasks that returned fewer than three samples are excluded from
the estimate rather than scored on what arrived (\cref{sec:metrics}), which is
why $n$ varies; failed requests are itemised in \cref{app:failures}. The
interval is a paired bootstrap over tasks (\cref{sec:paired}), not an
independent binomial one: every model answers the same items, and a binomial
interval here would be the estimator this paper spends \cref{sec:paired}
arguing against. Ordering is by \usable{}@3 and should be read as a banding,
not a ranking (\cref{sec:power}) --- the intervals overlap almost everywhere,
which is the point.}
\label{tab:main}
\end{table}

The apparent headline in \cref{tab:main} is the gap between $\pass{}@3$
(38.0--46.5\%) and $\usable{}@3$ (17.1--23.8\%): approximately half of all
test-passing generations fail at least one of the remaining two criteria, and
the ratio is stable across the entire price range.

\cref{sec:decomp} shows that this reading is wrong, or at least that it
attributes the gap to the wrong cause. We present it here as stated because it
is what the composite reports, and because constructing such a composite is a
natural thing to do. The point of the next section is that building one
without decomposing it was a mistake, and it was ours.

\subsection{Decomposing the gap}
\label{sec:decomp}

\usable{} conjoins three conditions. \cref{tab:decomp} attributes the drop
from \pass{}@3 to \usable{}@3 between the safety term and the equivalence
term, and reports alongside it the rate of ReDoS vulnerability
\emph{conditional on functional correctness}.

\begin{table}[htbp]
\centering
\footnotesize
\setlength{\tabcolsep}{4pt}
\input{tab_decomp}
\caption{Decomposition of the \pass{}$\to$\usable{} gap. ``C\&S'' is
correct-and-secure: correctness conjoined with safety, omitting the
equivalence term. The final column is the share of functionally correct
patterns that are ReDoS-vulnerable, and is the one column here computed
\emph{per sample} rather than @3: it is a rate over generations, and the @3
form would answer a different question. Denominators are \cref{tab:at1}'s
restricted ones --- tasks with fewer than three samples are excluded --- which
is why \texttt{kimi-k3} reads 46.5 here.}
\label{tab:decomp}
\end{table}

Safety costs \safetyloss{} percentage points on average. Equivalence costs
\equivloss{}. So the equivalence term supplies \textbf{\equivshare{}\%} of the
gap that made the composite look interesting.

That split is order-dependent, and the order we chose is the one that favours
the other term. Subtracting safety first credits it with every unsafe-correct
pattern, including those that would also have failed equivalence, while
equivalence keeps only the share no other term had already claimed. Reversing
the order can only move attribution toward equivalence. \equivshare{}\% is
therefore a \emph{lower bound} on the equivalence term's contribution under
any order-symmetric attribution, which is worth saying because the obvious
objection to a sequential decomposition is that it is order-dependent, and
here the dependence runs against us.

That would be unremarkable if the equivalence term were sound. It is not.
\cref{sec:goldstandard} shows, from an independent audit, that roughly 85\% of
what it counts against a model is reference-set error or an underspecified
prompt. Our composite was mostly reporting its own weakest component back to
us.

We did not see this coming. The decomposition was run late, after the audit
and after the first draft, because a reviewer of an earlier version asked how
our numbers compared to joint benchmarks in other domains. Answering that
required computing correct-and-secure without the equivalence term, and the
answer changed the paper.

Restricting attention to the two reference-independent criteria, the finding
is far more modest and far better supported: \textbf{\vulngivencorrect{}\% of functionally
correct regular expressions are ReDoS-vulnerable} (range 5.6--10.0\%).

The denominator matters, so: of the \correctsamples{} generations that
satisfied every test, \correctvulnsamples{} were simultaneously screened as
unsafe. That is \vulngivencorrect{}\%, a rate over \emph{samples} rather than
over tasks, which makes it an @1 quantity and the one directly comparable to
the single-sample figures in \cref{tab:srcompare}. At the task level the
corresponding count is 144 of 2{,}051 tasks on which some sample passed. About
one in fourteen either way, not one in two. (An earlier draft reported a count
of 135 here that we can no longer reproduce from our own released data, and
did not state which denominator the \vulngivencorrect{}\% used.)

$\exact{}@3$, string identity with the reference, runs from 2.0\% to 6.3\%.
String-comparison scoring would therefore misrank essentially every system,
which is what justifies the automata-theoretic treatment.

\subsection{Sample budget and the fragility of ordering}
\label{sec:atk}

Because $n=k=3$ renders the $\pass{}@k$ estimator degenerate
(\cref{sec:metrics}), we also report $@1$, which is the per-sample
success rate and the quantity most directly comparable to single-sample
protocols. \cref{tab:at1} gives both.

\begin{table}[htbp]
\centering
\small
\input{tab_at1}
\caption{Per-sample ($@1$) against any-of-three ($@3$) estimates, computed
from per-sample success counts. The final column gives the number of tasks
with fewer than three usable samples, which are excluded from that model's
$@3$ estimate. The $@3$ columns agree with \cref{tab:main} to the digit for
all eleven models, which is the check that matters: they are computed by
different routes over the same exclusion, and they used to disagree by up to
1.3 points. Reconciling them is how the estimator defect in
\cref{app:hazards} was found.}
\label{tab:at1}
\end{table}

The comparison is instructive beyond bookkeeping. At $@3$,
\texttt{kimi-k3} attains the highest \usable{} at 23.8\%. At $@1$ the leader
is \texttt{claude-opus-5} on 17.3\%, and \texttt{kimi-k3} drops to joint
third.
\emph{The induced ordering depends on the sample budget.} Systems differ not
only in per-attempt quality but in how much they benefit from additional
attempts, and a leaderboard that fixes $k$ silently fixes a weighting between
those two properties. This is a further reason to prefer banded reporting
(\cref{sec:power}), and a reason to state $k$ prominently whenever an @$k$
metric is ranked.

\subsection{ReDoS and the human baseline}
\label{sec:humanbaseline}

\cref{tab:vuln} applies the identical safety screen to the corpus's 450
human-written reference patterns. For comparability with the human set (one
pattern per task) we score one sample per task per model rather than the
any-of-three statistic of \cref{tab:main}.

\begin{table}[htbp]
\centering
\small
\input{tab_vuln}
\caption{ReDoS screening of model outputs against the corpus's own
human-written reference patterns, one pattern per task in all cases. The final
column is the exact McNemar test against the reference set on the same tasks;
see below for why an independent interval would be the wrong instrument here.}
\label{tab:vuln}
\end{table}

Every model screens safer than the reference set, and eight of the eleven do
so distinguishably. That second clause is the one that needs the work, and an
earlier draft did not do it: it read the ordering off eleven independent
binomial intervals over a pooled 4{,}941 patterns. That is the wrong estimator
twice over, and both errors are ones this paper criticises elsewhere. The
4{,}941 are eleven models answering the same 450 tasks, so task difficulty is
a shared random effect and a pooled binomial interval is too narrow
(\cref{sec:paired}). And each model is compared against the reference set
\emph{on the same tasks}, so the comparison is paired and an unpaired interval
discards the pairing --- with $n \approx 450$ a side and a three-point
difference, nothing would resolve.

McNemar's test is the paired procedure for this design, and we use the exact
binomial form because the discordant counts are small (47 to 65 per model).
Everything the model and the reference have in common cancels. Against
\texttt{qwen3.6-max-preview}, for instance, the reference set is vulnerable
where the model is safe on 39 tasks and safe where the model is vulnerable on
11 ($p = 0.0001$). Against \texttt{kimi-k3} the split is 39 to 26
($p = 0.14$), which does not resolve. The three models that do not separate
--- \texttt{kimi-k3}, \texttt{claude-opus-5} and \texttt{glm-5.2} --- are the
three with the highest vulnerability rates, which is what one would expect and
is not what an eleven-way ordering of point estimates would have told us.

So: models are safer than this corpus's answer key, on most of the field,
under the right test. Read alone, that invites a strong reading: ReDoS
vulnerability is endemic to human-written regular expressions, and models
merely reproduce it at a lower rate. We drafted that claim. It does not
survive the obvious check, which is that a benchmark answer key is not the
same population as working code.

\subsubsection{Six populations}
\label{sec:crosscorpus}

Screening a pattern requires no inference, so the human side of this
comparison can be extended to any corpus at no cost. We applied the identical
screen to five further populations: the gold answers of
KB13~\citep{kushman2013semantic}; the targets of
NL-RX~\citep{locascio2016neural}, generated from a grammar rather than
written by people and included as a control; and three corpora from the
artifact of \citet{davis2019linguafranca}, namely regular expressions
extracted from shipped packages across eight registries, regular expressions
posted to Stack Overflow, and the reusable patterns published on
\texttt{regexlib.com}. Samples are drawn at a fixed seed.

\begin{figure}[htbp]
\centering
\includegraphics[width=\linewidth]{fig_populations.pdf}
\caption{The same screen on every population, before and after controlling
for task mix. Bars are 95\% binomial intervals. The restriction is the same
rule on both sides --- keep only anchored \texttt{\^{}...\$} patterns --- and
it is applied to the models too, which an earlier draft did not do. What
survives it is not an ordering of humans against machines.}
\label{fig:populations}
\end{figure}

\begin{table}[htbp]
\centering
\small
\begin{tabular}{llrr}
\toprule
Population & Written to be & $n$ & vulnerable \\
\midrule
\multicolumn{4}{l}{\emph{All patterns}} \\
NL-RX-Synth (grammar-generated control) & --- & 5{,}840 & 35.2\% $\pm$ 1.2 \\
RegexLib, published for reuse & read & 3{,}446 & 17.2\% $\pm$ 1.3 \\
KB13 gold answers & read & 532 & 16.7\% $\pm$ 3.2 \\
Re(gEx$|$DoS)Eval gold answers & read & 755 & 13.1\% $\pm$ 2.4 \\
Stack Overflow posts & read & 5{,}000 & 5.9\% $\pm$ 0.7 \\
Production code & run & 5{,}000 & 4.9\% $\pm$ 0.6 \\
\addlinespace
\multicolumn{4}{l}{\emph{Anchored \texttt{\^{}...\$} only, controlling for task mix}} \\
RegexLib, published for reuse & read & 1{,}684 & 20.1\% $\pm$ 1.9 \\
Stack Overflow posts & read & 4{,}000 & 17.3\% $\pm$ 1.2 \\
Re(gEx$|$DoS)Eval gold answers & read & 538 & 13.4\% $\pm$ 2.9 \\
\textbf{This work, 11 models pooled} & \textbf{---} & 3{,}613 & \textbf{9.8\% [7.5, 12.4]} \\
\textbf{Production code} & \textbf{run} & 4{,}000 & \textbf{8.9\% $\pm$ 0.9} \\
\bottomrule
\end{tabular}

\caption{The identical ReDoS screen applied to six populations, plus this
paper's models under the identical restriction. Intervals are 95\% binomial,
which is the right form here and not in \cref{tab:main}: these populations
are genuinely independent --- different corpora, no shared items --- where
the models answer shared tasks and need a paired procedure. The anchored
block is the comparison to read: the unanchored figures are dominated by task
mix, which is why Stack Overflow moves from 5.9\% to 17.3\% between the two
blocks.}
\label{tab:crosscorpus}
\end{table}

Production code screens at 4.9\%, below every model we evaluated. The endemic
reading is therefore wrong, and we would have published it.

The raw comparison is confounded by task mix. This corpus is unusually rich
in validators (email, ISBN, hostname), the construction that backtracks,
whereas most regular expressions in the wild are short fragments with no
opportunity to. Restricting every population to anchored
\texttt{\^{}...\$} patterns gives the closest matched populations available,
and the ordering that emerges does not follow the axis we expected:

\begin{quote}
The dividing line is not human against machine. It is whether the pattern was
ever run. Every population written to be \emph{read} --- library entries at
20.1\%, forum answers at 17.3\%, benchmark reference answers at
\anchoredgold{}\% --- is ReDoS-prone. The one population that has been
\emph{executed} under real traffic sits at \anchoredprod{}\%, and the models
sit with it at \anchoredmodels{}\%.
\end{quote}

\paragraph{The model row is restricted the same way.} An earlier draft of this
table put the models into the anchored block at their \emph{unrestricted}
rate, taken straight from \cref{tab:vuln}. That compares an
anchoring-restricted human population against an unrestricted model one, which
is precisely the confound the restriction exists to remove, and a reviewer was
right to refuse it. The row now reports the models' rate on their own anchored
outputs --- the identical rule, applied to the pattern rather than to its
provenance --- which is \anchoredmodels{}\% $\pm$ \anchoredmodelsci{} over
\anchoredmodelsn{} patterns.

Correcting it moves the models \emph{toward} the answer key and away from
production code, by about eight tenths of a point, and the conclusion survives
anyway: \anchoredmodels{}\% against production code's \anchoredprod{}\% is not
a resolvable difference ($z = 1.27$, $p = \anchoredp{}$), while both sit far
below every population written to be read. We report the direction of the
correction because it cost us something.

Two stricter matchings are available here that are not available for the wild
populations, because we know which \emph{task} each model pattern answers.
Restricting to the 330 of 450 tasks whose reference is itself anchored, the
models screen at 10.6\% against the reference set's 13.9\%. Requiring both ---
an anchored output on an anchored task --- gives 11.4\% against 13.9\%, and at
that point the model-versus-reference difference is no longer resolvable
either. That is the honest ceiling on how strongly the reference comparison
can be stated, and it does not disturb the production comparison, which is a
population-level restriction applied identically on both sides.

This is a measurement of the mechanism rather than a conjecture about it. A
pattern published for reuse, posted in an answer, or written to key a
benchmark is authored once to communicate an idea, and nothing subsequently
happens to it. A pattern in a shipped package is executed, profiled, and
occasionally repaired; \citet{davis2018redos} document exactly that
remediation activity at ecosystem scale. Vulnerability tracks exposure to
execution, and answer keys have none.

A competing explanation produces the same ordering without any repair.
Packages under a lint rule such as \texttt{safe-regex} never contained the bad
pattern to begin with, so selection at authoring time would look exactly like
survival after it. The two are separable in principle --- repair leaves a
commit and prevention does not --- and we cannot separate them here: the
artifact we screen carries registry and module counts but no maintenance
history, so we can stratify by how widely a pattern is used and not by whether
its package is still maintained. Davis et al.'s documented remediation
supports the repair mechanism and does not rule out the other. We state the
finding as exposure-to-execution because that is what both mechanisms share,
and flag the decomposition as open.

It also explains why this corpus's answers are ReDoS-prone without appeal to
authorial carelessness. Re(gEx$|$DoS)Eval was assembled from real user posts,
and forum posts screen at 17.3\%. That the gold set is safer than its own
source population (\anchoredgold{}\%) suggests its authors curated, which is to their
credit and does not close the gap to executed code.

\paragraph{The control is the clearest case of the mechanism.} NL-RX-Synth
screens at 35.2\%, more than double the next population and seven times
production code. That is not the screen misbehaving on an odd construct
distribution, and the reasons say so: every one of the 2{,}057 vulnerable
patterns is flagged by the \emph{structural} pass, none by the empirical one,
and they divide into 994 with a quantifier wrapping a quantified group, 891
with two quantifiers in sequence over overlapping character sets, and 172 with
a quantifier over overlapping alternation. Representative targets are
\lstinline!.*(.*)([a-z]).*!, \lstinline!(dog.*[0-9].*)+! and
\lstinline!((dog)|(..*[AEIOUaeiou].*)){4,}!.

The generating grammar composes \texttt{.*}, \texttt{(\ldots)*} and
\texttt{(\ldots)\{n,\}} freely over overlapping character classes, and those
compositions \emph{are} the backtracking shapes. A grammar that samples
uniformly from a space of regular expressions has no reason not to, because
nothing downstream of it ever runs the output. Far from being an outlier that
impugns the screen, NL-RX-Synth is the purest instance of this section's
mechanism: it is the population with the least exposure to execution of any we
measured, and it is the most vulnerable. \cref{sec:replication} returns to
this, because StructuredRegex is grammar-built too.

\cref{sec:goldstandard} reaches the same conclusion from semantics, and this
section reaches it from safety without consulting the reference at all: the
reference set is the least reliable artifact in the evaluation. The revised
finding is weaker than the one we drafted, better supported, and preserves
the practical implication. Safety screening must be applied at the point of
use, because neither the generator, nor the reference author, nor the
benchmark performs it.

\subsubsection{Is the screen equally sensitive across these populations?}
\label{sec:calibration}

The ordering above is only a fact about the populations if the screen misses
vulnerable patterns at the same rate in all of them. If production patterns
are longer and structurally unlike RegexLib's showcase validators, a
differential miss rate could generate the entire result. ``\vuln{} is a lower
bound'' is an adequate caveat for one population and is not one for six, so we
measure it.

\paragraph{The screen does not see every dialect.} Every population is
screened by CPython's \texttt{re}, and patterns it will not compile are
dropped. That filter is not uniform, which is the part that matters: it
removes \proddroppedpct{}\% of the production pool (\proddropped{} of
\prodpool{}) against \sodroppedpct{}\% of Stack Overflow and
\libdroppedpct{}\% of RegexLib, and within production it runs from
\droprange{}. \cref{sec:threats} applies exactly this scrutiny to KB13 and
NL-RX, whose DSL \texttt{re} accepts as literals, and the corpus doing the
most work in this paper did not previously get it.

What is removed is \texttt{\textbackslash A\ldots\textbackslash z},
\texttt{\textbackslash G}, \texttt{\textbackslash Q\ldots\textbackslash E},
Perl's \texttt{\textbackslash x\{\ldots\}} and Unicode properties --- Perl and
Ruby syntax, and it is enriched for the anchored validator shapes the matched
comparison rests on: \dropanchoredpct{}\% of what is dropped is anchored
against \keptanchoredpct{}\% of what is kept. We had assumed this biased in
our favour and it does not
obviously do so: the read-to-be-read populations lose the larger share, and
losing vulnerable patterns from those would understate the very gap we report.
Assuming either direction is guessing, so we measure it. We translate what
can be translated without guessing (\texttt{runner/dialect.py}:
\texttt{\textbackslash A}$\to$\texttt{\^{}}, \texttt{\textbackslash
z}$\to$\texttt{\textbackslash Z}, \texttt{\textbackslash
Q\ldots\textbackslash E} to an escaped literal, named groups to Python's
spelling), decline what cannot (\texttt{\textbackslash G} anchors to the
previous match and has no Python equivalent; \texttt{\textbackslash
p\{\ldots\}} would have to be approximated), and report the recovered
population as a robustness row alongside the original.

\paragraph{Two registries cannot backtrack at all.} \linearregistries{} target
RE2 and the Rust \texttt{regex} crate, which are linear-time by construction.
ReDoS is not a property their patterns can have, whatever their shape, so
screening them as though they ran under a backtracking engine is a category
error. \linearpatterns{} patterns appear in one of the two; of those,
\linearonly{} appear in no backtracking registry at all, and those are the
ones the robustness run excludes. (A pattern can be published to several
registries, so this is a count of distinct patterns rather than the sum of the
two registry columns.)

\paragraph{Neither correction moves the number.} Normalisation recovers
\recovered{} of the \proddropped{} dropped production patterns
(\recoveredpct{}\%). What stays out is mostly not exotic syntax: the largest
group, \topunrecoveredn{} patterns, is malformed or truncated by the static
extraction that produced the corpus --- template fragments rather than
regular expressions --- followed by \texttt{\textbackslash G} and Unicode
properties, which we decline rather than approximate. Screening the recovered
population takes anchored
production from \anchoredprod{}\% to \robnorm{}\%; dropping the linear-engine
registries takes it to \robbt{}\%; doing both gives \robboth{}\%. The whole
range is \robrange{}\%, well inside the interval on any one of them, and the
models at \anchoredmodels{}\% remain indistinguishable from every variant.
The dialect filter and the engine mismatch are real defects in the
instrument, and they are not what produces the result.

\paragraph{Recall, measured per population.} Agreement between two static
approximations is not calibration, so we pair an independent detector with a
dynamic oracle. The detector is
\texttt{weideman-RegexStaticAnalysis}, from the artifact of
\citet{davis2019linguafranca} and the same one their ecosystem study used: a
different analysis, a different language, a different author. When it calls a
pattern vulnerable it emits an exploit string, and we then build that input at
growing pump counts and time \emph{CPython's own matcher} on it under a hard
timeout. A pattern is ground-truth vulnerable when matching either fails to
finish or grows measurably super-linearly in the input length. That is a fact
about the engine this paper screens with, not a second opinion.

\begin{table}[htbp]
\centering
\footnotesize
\setlength{\tabcolsep}{4pt}
\input{tab_calibration}
\caption{Screen sensitivity per population. \emph{Unanalysable} counts
patterns the independent detector skipped or timed out on, which are neither
evidence for nor against. \emph{Caught/confirmed} is our screen's hits over
the patterns the detector flagged \emph{and} a timing measurement then
confirmed blow up under CPython. Recall is measured against what the detector
finds, so it is a relative sensitivity rather than an absolute one --- which
is what the cross-population comparison needs, since the question is whether
the instrument is equally blind everywhere and not how blind it is. The
intervals are wide because dynamically confirmed vulnerabilities are scarce
in a sample of a few hundred; they are given so that the comparison is read
at the precision the counts support.}
\label{tab:calibration}
\end{table}

Recall runs \recallrange{}\% across the seven populations, and the intervals
overlap almost everywhere: dynamically confirmed vulnerabilities are scarce in
a sample of a few hundred, and KB13 yields three. So this does not resolve a
fine ordering of sensitivities, and we do not claim one.

What it does resolve is the direction, which is the part that mattered. If the
ordering were manufactured by the screen being blinder in production code than
in showcase validators, recall would have to be \emph{lower} in production.
It is not. The two populations where the screen is most sensitive are
\texttt{regexlib.com} at 91.7\% and Stack Overflow at 92.6\% --- the two with
the highest measured vulnerability --- while production code sits at 80.0\%
and this corpus's gold answers at 86.4\%. The differential runs the opposite
way to the one that would explain away the result.

Dividing each measured rate by the recall behind it gives a first-order
sensitivity check. On the anchored populations that turns 20.1\%, 17.3\%,
\anchoredgold{}\%, \anchoredmodels{}\% and \anchoredprod{}\% into
\corrlib{}\%, \corrso{}\%, \corrgold{}\%, \corrmodels{}\% and \corrprod{}\%.
Every rate rises, because the screen is a lower bound everywhere; the
\emph{ordering is unchanged}; and the read-against-run gap is not closed. We
would not put those corrected numbers in a table --- recall here is measured
against what one detector finds, so it is relative rather than absolute, and
the counts behind it are small --- but the check is the one a reader should
want, and it does not overturn anything.

Two limitations belong with it. The independent detector could not analyse
between 46 and 110 of each 240-pattern sample, and those are neither evidence
for nor against; ground truth therefore covers the patterns two tools can both
reason about, which is not all of them. And agreement between our screen and
the detector runs 53--74\%, far below either one's recall, because the
detector flags a good deal that does not blow up in CPython ---
\lstinline!^[a-z]+[a-z]*$! is quadratic in principle and instant in practice.
That gap is the reason this section uses a timing measurement as the arbiter
rather than a second opinion.

\paragraph{Relation to prior work, and a caveat on the split.}
\citet{siddiq2024redos} report that LLM-generated regexes skew toward
polynomial ReDoS. We do not reproduce this in our models: pooled, we observe
5.3\% exponential against 3.8\% polynomial. The cross-corpus run locates the
skew elsewhere. In the anchored production sample, polynomial exceeds
exponential 253 to 105, while this corpus's gold answers are near even at 34
to 38.

That comparison is weaker than it looks and we would rather say so than let it
be read as a finding. Our exponential-versus-polynomial split is only partly a
claim about growth order. The structural pass assigns it from the shape it
matched, which is principled; the empirical pass assigns it from \emph{which
probe length first timed out}, which is a threshold, not a measurement.
Patterns flagged only empirically therefore carry a degree label our
instrument did not really establish.

How much of the split that affects, we can say for one population. Among the
generations that satisfied every test, on both corpora, \emph{every} vulnerable
verdict came from the structural pass and none from the empirical one, so for
those the degree is shape-derived rather than threshold-derived. We have not
established the same for the reference sets or for the wild populations, where
patterns are longer and more of them fall outside the structural pass's
grammar. The subcategory counts are also small. We report the discrepancy with
prior work as an observation that wants a replication under a detector that
measures growth order directly, and draw nothing from it.

\subsection{Comparison against joint benchmarks in other domains}
\label{sec:crossbench}

Because \usable{} carries a semantic-equivalence conjunct that the joint
benchmarks of \cref{sec:jointprior} do not, comparing it directly to their
figures would understate us. We therefore also report
\emph{correct-and-secure}: correctness conjoined with safety, omitting the
equivalence term, which is the construction those benchmarks use.

\begin{table}[htbp]
\centering
\footnotesize
\setlength{\tabcolsep}{4pt}
\begin{tabular}{lrrrr}
\toprule
Benchmark & tasks & functional (\%) & joint (\%) & vuln.$\mid$correct (\%) \\
\midrule
This work (regex, ReDoS) & 450 & 42 & 39 & 7 \\
BaxBench \citep{vero2025baxbench} & 392 & 62 (best) & --- & $\approx$50 \\
SecureAgentBench \citep{chen2025secureagentbench} & 105 & --- & 15.2 (best), 9.2 (mean) & --- \\
DualGauge \citep{patir2025dualgauge} & 154 & $>$50 & $<$12 & --- \\
\bottomrule
\end{tabular}

\caption{Our correct-and-secure rate against joint correctness-security
results from other domains. Task types, vulnerability classes and evaluation
methods all differ substantially. The comparison is one of magnitude and of
the ratio between functional and joint success, not of models or difficulty.}
\label{tab:crossbench}
\end{table}

\textbf{We do not reproduce the pattern.} \citet{vero2025baxbench} report that roughly half of functionally correct
backends are exploitable. \citet{chen2025secureagentbench} report 15.2\%
correct-and-secure for their strongest agent, against substantially higher
functional correctness. \citet{patir2025dualgauge} report secure-pass@1 below
12\% while functional pass@1 exceeds 50\%. In each case the security criterion removes a large
majority of functionally correct output.

In our setting it removes \vulngivencorrect{}\%.

The one prior correct-and-secure figure on \emph{this} corpus is
\citet{siddiq2024regexeval}'s, over T5, Phi-1.5 and GPT-3.5-Turbo. We do not
put it in \cref{tab:crossbench}. Their generation, extraction and dialect
handling differ from ours in ways we did not control for, so a difference
between their number and ours would not be attributable to the two model
populations, which is the only comparison that would be interesting. Anyone
wanting that comparison should re-run their harness on the current models
rather than read it off the two papers.

We have several explanations for the divergence. One of them is testable on
this data and the test is in \cref{sec:replication}; the other two are not,
and we say which is which rather than listing all three as equally open:

\begin{enumerate}
\item \textbf{Attack surface.} A regular expression is a single expression
with one failure mode of interest. A backend application has many
independently exploitable components, so the probability that at least one is
vulnerable compounds with program size.
\item \textbf{Vulnerability class.} ReDoS is a structural property of the
pattern, plausibly well represented in training data as a recognised
anti-pattern. CWE-classified defects in application code are more diverse and
more context-dependent.
\item \textbf{Task difficulty ceiling.} Our functional pass rate is itself low
(38--47\%), so the correct subset may be biased toward simpler tasks whose
solutions have less structural room for catastrophic backtracking.
\end{enumerate}

The third is the one that would deflate the finding, and it is the one we can
test, because we have a second corpus that is twenty points easier. If the
correct subset is safe because it is simple, then StructuredRegex's correct
patterns should be simpler still and therefore safer. They are neither. On the
structural measures, StructuredRegex's correct patterns are \emph{longer}
(median 33 characters against 28) and carry \emph{more} quantifiers (median 4
against 3), while being twice as ReDoS-prone.
\cref{sec:replication} takes this further, because the direction of the
difference turns out to be informative about which shape is doing the damage.

The first two we cannot separate here. Both predict that the regex domain
looks safer than backend code, and neither predicts anything about the gap
between our two corpora, so nothing in this study distinguishes them.

So the claim that generalises is weaker than ``correct code is often
insecure,'' and more useful. \emph{The size of the correctness-to-security
penalty depends on the domain and has to be measured there.} Assuming it
transfers is what we did, and we were wrong. Worse, our composite showed a gap
of about the expected size for entirely the wrong reason
(\cref{sec:decomp}). An undecomposed joint metric will mislead you exactly
there, where its aggregate agrees with the literature and its components do
not.

\subsection{Replication on a second corpus}
\label{sec:replication}

Every figure above comes from one corpus, so the \vulngivencorrect{}\% is a property of
Re(gEx$|$DoS)Eval until shown otherwise. The two criteria that never consult
a reference can be computed anywhere that supplies test strings, and
StructuredRegex~\citep{ye2020structuredregex} supplies them: each instance
carries positive and negative example strings alongside its description. Its
own targets are in a prefix DSL and we never read them, so this run has no
$\dfaeq{}$ and no $\exact{}$ term, which suits the argument of this paper.

We collected one sample per task from the same eleven pinned models under the
same configuration, 622 of the 629 test instances (seven ship no positive or
negative strings, leaving $\pass{}@k$ undefined), for \$3.67.
\cref{tab:srcommon} reports the 513 tasks every model answered.
\cref{app:hazards} records why that number is not 622.

\begin{table}[htbp]
\centering
\small
\input{tab_sr_common}
\caption{StructuredRegex, one sample per task, on the 513 instances answered
by all eleven models. Percentages. Intervals are Wilson binomial, which is
appropriate here in a way it is not in \cref{tab:main}: this is one sample per
task, so there is no within-task resampling for a bootstrap to exploit. The
last column is the share of functionally correct patterns that are
ReDoS-vulnerable, the same quantity this paper reports as
\vulngivencorrect{}\% on Re(gEx$|$DoS)Eval.}
\label{tab:srcommon}
\end{table}

\begin{table}[htbp]
\centering
\small
\begin{tabular}{lcc}
\toprule
 & Re(gEx$|$DoS)Eval & StructuredRegex \\
\midrule
Tasks & 450 & 513 \\
Reference used in scoring     & yes ($\dfaeq{}$) & no \\
\midrule
$\pass{}@1$, range & 30.0--41.4\% & 51.7--66.1\% \\
$\vuln{}@1$, range & 7.4--10.7\% & 13.6--17.9\% \\
$\vuln{}\mid$correct, pooled & \textbf{7.4\%} & \textbf{16.5\%} \\
\bottomrule
\end{tabular}

\caption{The two reference-independent criteria across both corpora, at
$@1$ in each so the comparison is like for like.}
\label{tab:srcompare}
\end{table}

The qualitative result replicates and the quantitative one does not. Working
patterns are exploitable at a material rate on both corpora, which is the
finding. But the rate is \textbf{\srvulngivencorrect{}\%} here against
\vulngivencorrect{}\% on Re(gEx$|$DoS)Eval, and the obvious deflationary
reading, that harder problems produce worse patterns, is ruled out by the
direction of the difficulty gap: StructuredRegex is \emph{easier}, by roughly
twenty points of $\pass{}@1$.

\paragraph{Which shape, and why a grammar produces it.} The structure of the
correct patterns says where the extra vulnerability comes from, and it is not
where one would guess. StructuredRegex's correct patterns are longer (median
\srlen{} characters against \relen{}) and carry more quantifiers (median
\srquant{} against \requant{}), but they are \emph{less} nested: a quantifier
applied to a group --- the \texttt{(a+)+} shape everyone pictures when they
hear ReDoS --- appears in \srqgrouppct{}\% of them against \reqgrouppct{}\% on
Re(gEx$|$DoS)Eval, and their median group nesting depth is zero. The corpus
with fewer of the canonical exponential shapes is the corpus with twice the
vulnerability rate.

What it has instead is repetition in \emph{sequence}, and the screen's own
reasons say so. Of the vulnerable correct patterns on Re(gEx$|$DoS)Eval,
\renested{} are flagged for a quantifier wrapping a quantified group and
\readjacent{} for two quantifiers in sequence over overlapping character sets
--- \readjacentpct{}\% adjacent. On StructuredRegex the shares invert:
\srnested{} nested against \sradjacent{} adjacent, \sradjacentpct{}\%
adjacent. In absolute terms the nested count barely moves between corpora
while the adjacent count more than doubles. The extra vulnerability is
entirely the polynomial family: \texttt{\textbackslash s*\textbackslash s*},
\texttt{[a-z]+[a-z0-9]*}.

Which is what its targets are built out of. A grammar over repetition,
optionality and concatenation produces adjacency by construction, and
\cref{sec:crosscorpus} found the identical signature in NL-RX-Synth, the
other grammar-built corpus here: 891 of its 2{,}057 vulnerable targets are
flagged for that same shape, and the population screens at 35.2\%. Two
grammar-generated corpora, built by different authors for different purposes,
over-represent the same family.

We think the reading is this. A grammar samples compositions uniformly and has
no reason to avoid adjacency, because nothing downstream of it runs the
output; a human writing a validator writes one nested quantifier and stops.
So the correctness-to-security penalty depends not only on the domain but on
how the corpus's targets were \emph{authored}, which is a sharper and more
uncomfortable claim than the one we set out with.

\cref{sec:crossbench} argues from BaxBench and SecureAgentBench that the
correctness-to-security penalty is domain-dependent. That argument compares
across languages and vulnerability classes, and a reader is entitled to
discount it. The same conclusion now holds between two corpora of the same
task, in the same language, under the same screen and the same eleven
models, with the effect size differing by a factor of 2.2. We regard this as
the stronger form of the claim, and it cuts against our own headline number.

The ordering of models is also not preserved. \texttt{claude-sonnet-5} leads
here and sits mid-table on Re(gEx$|$DoS)Eval; \texttt{kimi-k3} leads there
and sits eighth here. \cref{sec:power} declines to publish a ranking on
single-corpus evidence, and this is the direct test of that caution.

\subsection{Cost}

\begin{table}[htbp]
\centering
\small
\input{tab_cost}
\caption{Cost per \emph{request}, measured from per-request billing returned
by the serving layer rather than estimated from list prices. A task costs
three of these at $n=3$: \texttt{claude-opus-5} at $2514.2 \times 10^{-6}$
over 1{,}350 requests gives the \$3.39 total.}
\label{tab:cost}
\end{table}

The extremes differ by a factor of 98 in price and 1.1 percentage points in
\usable{}@3 --- and in the direction that favours the cheaper model, which
scores 19.8\% against 20.8\%. That difference is well inside what we can
resolve (\cref{sec:paired}), which is to say the two are indistinguishable.
This observation is about this task only. It is not a general claim about
model capability.

\subsection{Refusals}

A content filter blocked \texttt{claude-opus-5} on 29 requests, 2.1\% of its
calls, spread across five tasks. Four have some security flavour: strings that
exclude a single quotation mark, a six-character alphanumeric password, hex
byte sequences, and SQL update and insert statements. The fifth is
\emph{``Matches a file extention.''} No other model refused anything.

The refusals were not deterministic. Several tasks came back refused on one
sample and answered on the other two, under identical settings. You can only
see that at $k>1$. A single-sample protocol would have logged a flat failure
and moved on. We count refusals as their own failure category, not as wrong
answers (\cref{app:failures}).

\subsection{Inference-time compute}

On the corpus's simplest task (\emph{``Matches exactly 1 numeric digit
(0-9).''}), \texttt{qwen3.6-max-preview} emitted 1{,}571 hidden reasoning
tokens before answering \texttt{\^{}[0-9]\$}, at a cost of \$0.0098. With
reasoning disabled it produced the identical answer in 10 tokens at
$1/100$th the cost. Asking that model for \emph{low} reasoning effort produced \emph{2{,}375}
reasoning tokens, more than the default. Effort controls are not reliably
monotone.

A reasoning-enabled comparison on 12 tasks across five models yielded
per-request costs $15.7\times$ higher and no interpretable accuracy signal at
that sample size ($\pm14$ percentage points per cell). We report the cost
result and explicitly decline to report an accuracy conclusion.

\section{Measurement Validity}
\label{sec:validity}

\subsection{The reference standard contains errors}
\label{sec:goldstandard}

A benchmark's answer key is a set of labels, and labels have an error rate.
\citet{northcutt2021labelerrors} audit ten widely used test sets and find a
mean 3.3\% error rate, enough in several cases to reverse the ordering of
models the benchmark was being used to compare. The corpus here differs in a
way that helps: a label is a regular expression, so a disputed label can be
settled by exhibiting a string on which the reference and the description
disagree, rather than by a second opinion. Every adjudication below rests on
such a string.

First we ruled out a loading fault. 449 of 450 reference patterns satisfy
their own tests, so the corpus semantics and dialect are configured correctly.
The real defect is quieter. A reference can pass its own tests and still fail
to denote what its description says, whenever the tests never exercise the
difference.

We drew a seeded random sample of 60 cases in which a model \emph{passed
every test} yet received a \textsc{diff} verdict, and adjudicated 14 in
detail. The 14 are \emph{the first 14 in seed order whose witness string is
ASCII}. That rule matters for the arithmetic below, so it is stated rather
than left to be recovered from the released sample: the non-ASCII cases are a
known and separate artifact, and adjudicating them would have been
adjudicating Python's Unicode defaults rather than either party's regular
expression.

\begin{table}[htbp]
\centering
\small
\begin{tabular}{lrr}
\toprule
Adjudicated cause & count & share \\
\midrule
Reference pattern is incorrect & 5 & 36\% \\
Description underspecifies the disputed property & 6 & 43\% \\
Model output is incorrect & 3 & 21\% \\
\bottomrule
\end{tabular}
\caption{Adjudication of 14 sampled \textsc{diff} verdicts on test-passing
outputs. Per-case reasoning is released.}
\label{tab:adjudication}
\end{table}

Separately, 19 of the 60 sampled disagreements (32\%) differ only on
non-ASCII input. Python's \texttt{\textbackslash d} is Unicode-aware and
\texttt{[0-9]} is not, so the two are formally distinct languages and
practically the same thing.

Because the adjudicated 14 were drawn from the ASCII-witness cases only, none
of them is among those 19, and the two effects compose without
double-counting: the model is at fault in 3 of 14 ASCII-witness cases, and
41 of the 60 sampled cases have an ASCII witness, so the model is clearly at
fault in $(41/60) \times (3/14) = 14.6\%$ of \textsc{diff} verdicts. Call it
roughly \textbf{15\%}. The verdict labels in \cref{tab:adjudication} map to
the released adjudications as \texttt{model\_right} $\to$ reference
incorrect, \texttt{ambiguous} $\to$ description underspecifies,
\texttt{gold\_right} $\to$ model incorrect.

Representative cases:

\begin{itemize}
\item \textbf{Description:} \emph{``It just accepts only positive numbers.''}
The reference \lstinline!^\d+([.,]?\d+)?$! admits \texttt{0}. The model
excluded zero and was scored incorrect.

\item \textbf{Description:} \emph{``A very simple ISBN validation
expression---it just checks for a 10 digit number.''} The reference is
\lstinline!^\d{9}[\d|X]$!. The character class contains digit, \emph{pipe}
and \texttt{X}: an alternation operator written inside a class, where it
denotes a literal. The reference therefore accepts \texttt{000000000|}, which
is the witness separating it from the model's \lstinline!^\d{9}[\dX]$!.

\item \textbf{Description:} \emph{``\ldots make sure commas are in the rite
place (if present).''} The reference
\begin{lstlisting}
^\$?\d{1,3}(,?\d{3})*(\.\d{1,2})?$
\end{lstlisting}
makes the comma optional, admitting \texttt{0,000000} and negating the stated
purpose. The model enforced comma placement and was scored incorrect.

\item \textbf{Underspecification:} \emph{``Matches any single upper- or
lower-case letter.''} The reference \lstinline!^[a-zA-Z]$! requires a
whole-string match. A candidate \lstinline![A-Za-z]! requires only
occurrence. The description does not distinguish the two. We tested whether anchoring alone
explains failures generally: it accounts for approximately 6\%, so the effect
is diffuse rather than a single correctable convention.
\end{itemize}

\paragraph{Consequence.} \dfaeq{} and \usable{} are lower bounds on model
correctness by an unmodelled margin. \pass{} and \vuln{} are unaffected,
since neither consults the reference.

\paragraph{Limitation of this analysis.} Fourteen cases, judged by us, on our
own benchmark. That is a weak basis for a precise number and we do not want it
cited as one. We release the reasoning for each case, and the unadjudicated
cases keep an empty verdict field so someone else can fill them in. Anyone
quoting 15\% should re-derive it from a larger sample with annotators who have
no stake in the answer. We report it because the direction is not in doubt,
and because the direction is what changes how the composite should be read.

\subsection{The composite rewards putting an answer beyond checking}
\label{sec:undec}

\dfaeq{} counts an \textsc{undec} verdict as a failure. \usable{} does the
opposite: its third conjunct excludes only a \emph{proven} \textsc{diff}, so a
comparison the engine cannot answer is scored as a pass. Each convention is
defensible alone --- one is a lower bound that cannot flatter, the other
isolates the model from the engine's reach --- but together they create an
incentive. Equivalence is undecidable exactly when a pattern uses
backreferences, so a model that reaches for one is scored \emph{more} usable
for having put its answer beyond checking. Between 78 and 133 tasks per model
land there, a spread of 55, which is large enough to matter.

The obvious test is to correlate a model's \textsc{undec} count against its
\usable{}@3. That returns nothing (Pearson $-0.13$), and it returns nothing
misleadingly: undecidability credit and equivalence skill push the composite
in opposite directions, and a correlation on the sum nets them out. So we
decompose instead. \cref{tab:undec} recomputes \usable{}@3 under the stricter
rule that the verdict must be \emph{proven} \textsc{eq}.

\begin{table}[htbp]
\centering
\small
\begin{tabular}{lrrrrr}
\toprule
Model & undec. & \usable{}@3 & proven-\textsc{eq} only & tasks & share of \usable{} (\%) \\
\midrule
\texttt{kimi-k3} & 78 & 23.8 & 13.2 & 47 & 44.8 \\
\texttt{qwen3.6-max-preview} & 94 & 21.6 & 12.7 & 40 & 41.2 \\
\texttt{gpt-5.6-sol} & 133 & 20.9 & 8.7 & 55 & 58.5 \\
\texttt{claude-opus-5} & 104 & 20.8 & 10.6 & 44 & 48.9 \\
\texttt{deepseek-v4-flash-0731} & 93 & 20.0 & 10.9 & 41 & 45.6 \\
\texttt{qwen3.6-plus} & 99 & 19.8 & 10.2 & 43 & 48.3 \\
\texttt{glm-5.2} & 86 & 18.7 & 9.6 & 41 & 48.8 \\
\texttt{gpt-5.6-terra} & 100 & 18.7 & 9.8 & 40 & 47.6 \\
\texttt{gpt-5.6-luna} & 112 & 18.5 & 8.5 & 45 & 54.2 \\
\texttt{claude-sonnet-5} & 97 & 18.0 & 9.8 & 37 & 45.7 \\
\texttt{gemini-3.1-flash-lite} & 95 & 17.1 & 8.7 & 38 & 49.4 \\
\midrule
\textit{pooled} & --- & --- & --- & 471 & \textbf{48.4} \\
\bottomrule
\end{tabular}

\caption{What the undecidability convention is worth. ``Proven-\textsc{eq}
only'' requires a demonstrated equivalence rather than the absence of a
demonstrated difference; the final columns are the tasks that separate the
two.}
\label{tab:undec}
\end{table}

\undecshare{}\% of all \usable{}@3 credit --- \undecsupported{} of
\undecusable{} task-credits --- rests on a verdict the engine could not reach
rather than on demonstrated equivalence, and the share tracks \textsc{undec}
count closely across models ($r = +0.82$) even though the composite itself
does not. Removing the credit reorders the field. \texttt{gpt-5.6-sol} has the
most undecidable comparisons of any model at 133; it has the third-highest
\usable{}@3 and one of the three lowest proven-equivalence rates, and it falls
seven places.

We are not proposing the stricter rule --- it would charge a model for a
theorem, which is the objection to \dfaeq{} --- and we are not claiming any
model games this deliberately. The point is narrower and it is a third
independent reason to distrust the composite: nearly half of what
\usable{} certifies, it certifies by default.

\subsection{Paired inference}
\label{sec:paired}

Our first analysis compared systems with independent binomial confidence
intervals. That is the wrong estimator for this design. All eleven models
answer the same items, so task difficulty is a shared random effect. Treating
the systems as independent samples charges every one of them separately for
the same variation.

A paired bootstrap fixes this, and the fix is not ours: the procedure is
\citet{bergkirkpatrick2012empirical}, following \citet{koehn2004bootstrap},
and \citet{dror2018hitchhiker} recommend it for exactly this design. Resample
tasks with replacement~\citep{efron1994bootstrap}, score every
model on each resample, and form the per-pair differences. Task difficulty
cancels, and only disagreement between the two systems contributes. With $B = 10{,}000$
resamples and a fixed seed:

\begin{table}[htbp]
\centering
\small
\begin{tabular}{lr}
\toprule
Procedure & pairs resolved at 95\% (of 55) \\
\midrule
Independent binomial intervals & 1 \\
Paired bootstrap, \usable{}@3 & 9 \\
Paired bootstrap, \pass{}@3 & 15 \\
Paired bootstrap, \vuln{}@3 & 15 \\
\bottomrule
\end{tabular}
\caption{Resolvable pairwise comparisons under independent versus paired
inference on identical data.}
\label{tab:paired}
\end{table}

The correction is not marginal. It changes the number of defensible
comparisons ninefold. We flag it because independent
intervals look like the default in adjacent evaluation work, and the design
they are applied to, all systems answering all items, is close to universal.

Even corrected, the middle of the table does not separate. For example
\texttt{qwen3.6-max-preview} versus \texttt{gpt-5.6-sol} yields
$+0.7\%$ with a 95\% interval of $[-2.9\%, +4.3\%]$.

\subsection{Discriminative power of the corpus}
\label{sec:power}

There is a structural reason behind this. On \usable{}@3, \textbf{62\% of
tasks produce the same outcome for all eleven models}. On \pass{}@3 it is
56\%, on \vuln{}@3, 77\%. Only 167 of 450 tasks carry any
discriminative signal at all.

For calibration: resolving a 3-percentage-point difference at $p \approx 0.2$
with independent intervals requires roughly 2{,}800 items---more than the
full 762-task corpus contains. Paired inference reduces this requirement
substantially but cannot manufacture signal from items on which every system
agrees. Benchmark construction for system comparison should therefore
optimise for discriminative yield, not corpus size.

\section{Threats to Validity}
\label{sec:threats}

\paragraph{Contamination.} The corpus predates every model we evaluated and
is public, as are the forum posts it was built from. Some of what we measured
may be memorisation. \citet{sainz2023contamination} argue that this has to be
measured per benchmark rather than assumed negligible, and we cannot measure
it: we have no held-out private task set, so we cannot put a bound on it.
This is the largest unaddressed threat in the paper.

\paragraph{Corpus-local effect sizes.} \cref{sec:replication} replicates
$\pass{}$ and $\vuln{}$ on a second corpus and finds the vulnerability rate
among correct patterns differs by a factor of 2.2. Every effect size in this
paper should therefore be read as a property of Re(gEx$|$DoS)Eval rather than
of natural-language-to-regex generation. $\dfaeq{}$ and $\exact{}$ remain
single-corpus: replicating those needs a second corpus with a reference in
standard syntax, and neither KB13, NL-RX nor StructuredRegex supplies one. A further hazard we avoided
rather than solved: KB13 and NL-RX are written in a DSL where \texttt{\&} is
intersection and \texttt{\~{}} is negation, both of which Python's
\texttt{re} accepts silently as literals. We excluded the affected patterns rather than mistranslate them: NL-RX-Synth
goes from 10{,}000 rows to 9{,}648 distinct targets to the 5{,}840 screened,
and KB13 from 824 rows to 732 distinct to 532. (An earlier draft gave the
exclusion counts over raw rows and the pool sizes over distinct ones, which
did not reconcile.)

\paragraph{Single prompt.} Every result here is conditional on one
unoptimised prompt. \cref{sec:goldstandard} shows how often the task
descriptions leave the anchoring convention unstated, so sensitivity to
phrasing is plausibly large. We did not measure it.

\paragraph{Sampling control.} We do not set temperature
(\cref{sec:sampling}), so it varies by provider default. We checked that
outputs still vary across samples. We did not equalise the sampling.

\paragraph{Estimator degeneracy.} At $n=k=3$ the $\pass{}@k$ estimator
reduces to any-of-3.

\paragraph{Coverage asymmetry.} Nine models cover all 450 tasks.
\texttt{kimi-k3} covers 447, having exhausted our budget mid-run, and
\texttt{claude-opus-5} covers 444, losing five tasks to content-filter
refusals and one to replies that were nothing but a code fence
(\cref{app:failures}). Tasks that came back
with fewer than $k$ samples are excluded from that model's @$k$ estimate
rather than scored on what arrived (\cref{sec:metrics}), so the denominators
in \cref{tab:main} are 432--450.

\paragraph{Screening, not proof.} \vuln{} is a lower bound
(\cref{sec:metrics}).

\paragraph{Single corpus.} We do not test whether the observed banding
survives corpus change.

\paragraph{Self-adjudication.} See \cref{sec:goldstandard}.

\section{Recommendations for Evaluation Design}

\begin{enumerate}
\item \textbf{Partition metrics by reference dependence, and decompose any
composite that spans the partition.} Metrics evaluated by execution against
inputs (\pass{}, \vuln{}) are unaffected by reference-set defects. Metrics
evaluated against a human answer key inherit its error rate. A composite
conjoining both can report a large gap that is almost entirely attributable to
the unreliable term---in our case \equivshare{}\% of it (\cref{sec:decomp})---while
appearing to be a statement about the reliable ones.

\item \textbf{Sample and adjudicate disagreements.} It costs hours. In our
case it changed the conclusion.

\item \textbf{Use paired inference when all systems answer all items.}
\cref{tab:paired}.

\item \textbf{Report discriminative yield}, not only corpus size
(\cref{sec:power}).

\item \textbf{Pin the serving provider and verify post hoc.} Record what
served each request, not only what was requested (\cref{sec:serving}).

\item \textbf{Instrument the scorer with controls} that must pass and must
fail (\cref{sec:metrics}).

\item \textbf{Release raw generations.} Anyone can then overturn a scoring
decision without buying the inference again.

\item \textbf{Treat an effect size as corpus-local until replicated.} Our
reference-independent gap is \vulngivencorrect{}\% on one corpus and \srvulngivencorrect{}\% on another under
an identical screen (\cref{sec:replication}). Reference-independence buys
freedom from answer-key error. It does not buy generality.

\item \textbf{Check that a partial run is not a biased run.} When a model
loses tasks to infrastructure, score the remaining models on the tasks it
lost. Ours were 8.1 points harder (\cref{app:hazards}), which would have
flattered the affected model rather than penalised it. Check separately that
the \emph{estimator} is defined on the partial data: ours was not, and
credited every short task to whichever model lost it (\cref{app:hazards}).

\item \textbf{Generate every table from committed data, and let the build
fail.} The reviewer of an earlier draft of this paper found nine arithmetic
inconsistencies. Every one of them was in a table we maintained by hand or in
a sentence quoting one, and none was in a table generated from the scores.
The fix is mechanical rather than editorial: emit all of them, keep the
numbers the prose states in macros generated alongside, and have continuous
integration rebuild and diff.

\item \textbf{Calibrate a screen per population before comparing
populations.} A detector's false-negative rate is a property of the patterns
it is pointed at. Reporting it as a single caveat is adequate for one
population and lets a differential miss rate masquerade as a finding across
several (\cref{sec:calibration}).
\end{enumerate}

\section{Conclusion}

We set out to measure a gap between correct and shippable regular
expressions. We built a composite metric, it showed a large gap, and
decomposition showed that \equivshare{}\% of that gap came from the one term we had
already shown to be unreliable. What survives is much smaller. \vulngivencorrect{}\% of
functionally correct patterns carry a ReDoS vulnerability, against roughly
50\% of functionally correct backends in comparable work. The
correctness-to-security penalty depends on the domain, and in ours it is
small.

Two findings survive that correction intact, both resting only on criteria
that never consult a human answer key. Models are safer than the
human-written reference patterns against which they are scored --- on eight of
eleven under a paired test, not on all eleven, which is what we would have
claimed from the point estimates --- and indistinguishable from the regular
expressions in shipped packages once both are restricted the same way. Across
six populations, ReDoS prevalence tracks whether a pattern is ever executed
rather than whether a human or a model wrote it, which locates the anomaly in
the answer key rather than in generation or in human practice. And the entire
eleven-model field lies within seven percentage points across a $98\times$
range in inference cost.

A third check went the same way. The \vulngivencorrect{}\% that survived the first two
corrections is \srvulngivencorrect{}\% on a second corpus under the same screen and the same
models (\cref{sec:replication}), so the number we were left defending is
corpus-local too.

All three corrections ran the same way. A composite agreed with the
literature and turned out to be dominated by its weakest term; a human
baseline on one corpus supported a strong claim that a second population
refuted; a reference-independent rate proved to be a fact about the corpus.
In each case the check was cheap and we nearly skipped it. The recurring
lesson is not about regular expressions. It is that a number which agrees
with what you expected is the one most worth spending an afternoon
attacking.

The methodological result is the one we did not plan on. A metric that
compares against a human answer key inherits that key's error rate, and a
composite built on such a term can produce a finding that is large, plausible
and mostly spurious. We only found this by reading our own output after
collecting it. Fourteen cases were enough to change what we were willing to
publish off a \$10.95 run. We would read them earlier next time.

\section*{Acknowledgements}

Scoring is performed by \texttt{regexbench}~\citep{regexbench}, prior work by
the present authors, used here as a pinned dependency and not contributed by
this paper. The corpus is Re(gEx$|$DoS)Eval~\citep{siddiq2024regexeval}, used under its
original terms and not redistributed.

\bibliographystyle{plainnat}
\bibliography{refs}

\appendix

\section{Failure taxonomy}
\label{app:failures}

Of 14{,}850 requests, 54 (0.36\%) failed and are reported rather than
excluded.

\begin{table}[htbp]
\centering
\small
\begin{tabular}{llr}
\toprule
Model & Cause & Count \\
\midrule
\texttt{claude-opus-5} & content-filter refusal & 29 \\
\texttt{claude-opus-5} & code fence, no pattern inside & 7 \\
\texttt{kimi-k3} & account spending limit reached mid-run & 11 \\
\texttt{kimi-k3} & response without resolved provider & 3 \\
\texttt{kimi-k3} & code fence, no pattern inside & 2 \\
\texttt{gpt-5.6-sol} & response without resolved provider & 1 \\
\texttt{gpt-5.6-luna} & code fence, no pattern inside & 1 \\
\midrule
& \textit{total} & \textit{54} \\
\bottomrule
\end{tabular}
\caption{Request failures by cause. Responses lacking a resolved provider are
discarded rather than scored, as provenance is a precondition for
reproducibility. The code-fence rows are responses the API returned as
successful, from which extraction yields the empty string: nine emitted a
fence with nothing inside it beyond an optional language tag, and one wrote a
pattern followed by an unmatched closing fence, which leaves the last fenced
block empty. All ten are short replies of 2--18 completion tokens; none is a
truncation against the 400-token cap. An earlier draft of this table omitted
them and accounted for only 44 of the 54.}
\end{table}

Empty completions are classified by reported \texttt{finish\_reason} into
refusal, token-budget exhaustion, and unexplained, rather than being scored
as empty patterns---which would be indistinguishable from a maximally poor
answer.

\section{Operational hazards}
\label{app:hazards}

We document configuration hazards encountered during this study, each of
which produced plausible-looking but invalid data. We believe this list has
independent value.

\begin{enumerate}
\item \textbf{Parameter combinations that are silently unroutable.} Setting
\texttt{temperature} together with \texttt{require\_parameters=true} yields a
hard routing failure on 6 of 11 models. Our first pilot failed all 396 calls.

\item \textbf{Systems that cannot satisfy the experimental condition.} Two
Gemini variants return \emph{``Reasoning is mandatory for this endpoint''}
and cannot be evaluated with reasoning disabled. Undetected, this would have
placed one reasoning-enabled system in an otherwise reasoning-disabled
comparison.

\item \textbf{Empty completions from budget exhaustion.} At a 200-token cap,
reasoning models consumed the entire budget on hidden reasoning and returned
no content.

\item \textbf{Resumption across configuration changes.} Our collection is
resumable by (model, task, sample). An aborted run at one token cap left
truncated rows that a resumed run at a different cap treated as complete,
interleaving two configurations indistinguishably. Rows now carry a
configuration fingerprint and resumption refuses on mismatch.

\item \textbf{Sample collapse in scoring.} An early scorer retained only the
final sample per task, which would have rendered all $@3$ statistics
disguised $@1$ statistics with no visible symptom.

\item \textbf{An estimator guard that is sound only on its own domain.} The
$\pass{}@k$ estimator is usually implemented with the shortcut ``if
$n - c < k$, return 1'': if fewer than $k$ samples failed, every choice of $k$
must contain a success. That is correct for $n \geq k$, which is the only case
the estimator is defined on, and it is silently wrong below it. A task that
returned one sample satisfies $n - c \leq n < k$ whatever happened, so it
scores a full 1.0 on every metric --- $\mathrm{pass\_at\_k}(1, 0, 3) = 1$ ---
and a task nobody answered correctly is counted as answered correctly.

We found this late, from a discrepancy of 1.3 points between two of our own
tables, and it had been inflating exactly the two models that lost the most
samples, which were the two at the top of the table. Correcting it moves
\texttt{claude-opus-5} from second to fourth on \usable{}@3 and makes
\texttt{kimi-k3} rather than \texttt{claude-opus-5} the highest $\pass{}@3$.
No claim in this paper turns on it --- we publish no ranking, and
\cref{sec:power} already says the band does not separate --- but
\cref{tab:main} is ordered by the affected column, and a reader comparing two
tables would have had no way to tell which was right.

The general form is worth naming: a fast path justified by a precondition,
in code that is also called when the precondition does not hold. The failure
is silent, it is biased toward the systems with the most missing data, and
nothing in the output looks wrong.
\end{enumerate}

\paragraph{A content filter can silently bias a subset.}
Collecting StructuredRegex, Anthropic's content filter rejected 182 of 622
descriptions for \texttt{claude-opus-5}, all of them benign (``A string that
starts with one or more digits and optionally ends with `NU' or `DG'.'').
A probe suggested the filter was stochastic: 8 of 20 rejected prompts
succeeded when resent unchanged. Five retry rounds recovered 98, with the
per-round yield collapsing (51, 31, 9, 2, 5), which shows a stochastic
component over a hard core the filter rejects every time. That leaves 84
prompts the filter rejected on every attempt. A further 25 responses came back
successful but yielded no pattern, and are not the filter's doing: 24 are a
bare code fence of one or two tokens, the same failure mode as
\cref{app:failures}'s, and one is a genuine truncation, having spent the whole
400-token budget inside a \texttt{<thinking>} block. Together
$84 + 25 = 109$, and final coverage was $622 - 109 = 513$.

The retry effort was not uniform across models, because nothing else needed
it: only \texttt{claude-opus-5} was blocked, so only \texttt{claude-opus-5}
received five extra rounds. The recovered 98 are therefore conditioned on
non-refusal in a way the other ten models' answers are not. We do not think
this biases regex quality --- the filter keys on the description, which the
model has not yet answered --- but it is differential effort and the
common-subset comparison below is the reason it does not matter for the
reported numbers.

The recovery is not the interesting part. Scoring the other ten models on the
109 blocked tasks against the 513 kept shows the blocked tasks are
$8.1$ points harder on average (sd $3.3$, negative for all ten). The
surviving subset therefore \emph{flattered} the affected model, at the same
time as counting the blocked rows as failures \emph{penalised} it. Either
error alone is visible; together they produce a plausible number that means
nothing. We report the common 513-task subset for all eleven models
(\cref{tab:srcommon}) and release \texttt{runner/filter\_bias\_check.py},
which needs no API access.

\section{Reproduction}

\begin{lstlisting}[language=bash]
git clone https://github.com/plicara/regexeval-2026
cd regexeval-2026
make setup      # pinned scorer + corpus download
make score RUN=sweep
\end{lstlisting}

Scoring reads only the committed generations. No API access, no
expenditure. We verified reproduction from a clean checkout with fresh
dependency installation and corpus download, obtaining bit-identical metrics.
A reduced form of this check executes in continuous integration on every
commit.

\end{document}
